\begin{definition}[Iterator-Generated Function]
\label{def:iterator-generated-function}
Let $(P,S,1)$ be a Peano system, and let $(W,c,g)$ be iterator data for it.
The unique function $f:P\to W$ satisfying
\[
f(1)=c
\qquad\text{and}\qquad
\forall n\in P\;(f(S(n))=g(f(n)))
\]
is called the \emph{iterator-generated function determined by $(W,c,g)$}.
\end{definition}
\end{definitionbox}
\LeanFormalizes{def:iterator-generated-function}{lra-lean}{LRA.VolumeII.PeanoSystems.TexSkeleton}{IteratorGeneratedFunction}{pending}

\begin{remark*}[Standard quantified statement]
\[
f \text{ is the iterator-generated function determined by } (W,c,g)
\Longleftrightarrow
\bigl(f:P\to W\land f(1)=c
\land \forall n\in P\;(f(S(n))=g(f(n)))\bigr).
\]
\end{remark*}

\begin{remark*}[Predicate reading]
\[
\operatorname{IteratorGeneratedFunction}(f,(P,S,1),W,c,g)
\Longleftrightarrow
\operatorname{IteratorFunction}(f,(P,S,1),W,c,g)
\land
\exists! h:P\to W\;
\operatorname{IteratorFunction}(h,(P,S,1),W,c,g).
\]
\end{remark*}

\begin{remark*}[Negated quantified statement]
\[
\begin{aligned}
f \text{ is not the iterator-generated function determined by } (W,c,g)
\Longleftrightarrow\;&
f \text{ is not a function } P\to W\\
&\lor f(1)\neq c
\lor \exists n\in P\;(f(S(n))\neq g(f(n))).
\end{aligned}
\]
\end{remark*}

\begin{remark*}[Failure modes]
\begin{description}
  \item[Exposition.]
  A candidate fails to be the iterator-generated function if it is not a
  function from $P$ to $W$, if it has the wrong base value, or if it violates
  the successor clause.
  \item[Invalid candidate.]
  \[
  \neg\operatorname{IteratorGeneratedFunction}(f,(P,S,1),W,c,g).
  \]
  \[
  \begin{aligned}
  f \text{ is not the iterator-generated function determined by } (W,c,g)
  \Longleftrightarrow\;&
  f \text{ is not a function } P\to W\\
  &\lor f(1)\neq c
  \lor \exists n\in P\;(f(S(n))\neq g(f(n))).
  \end{aligned}
  \]
\end{description}
\end{remark*}

\begin{remark*}[Interpretation]
This definition packages the unique function supplied by the Peano Iterator
Theorem as a named object available for later constructions.
\end{remark*}

\begin{dependencies}
\begin{itemize}
  \item \hyperref[thm:peano-iterator-theorem]{Peano Iterator Theorem}
\end{itemize}
\end{dependencies}

